% ============================================================
% content/answers-short.tex —— 2022 湖南卷物理 · 简明答案速查
% ============================================================

\begin{gknotice}
  选择题答案：1--C，2--A，3--D，4--B，5--A，6--B，7--AC，8--CD，9--BC，10--BD
\end{gknotice}

\section*{一、选择题}

\noindent
\begin{tabular}{@{}lcccccccccc@{}}
\toprule
题号 & 1 & 2 & 3 & 4 & 5 & 6 & 7 & 8 & 9 & 10\\
\midrule
答案 & C & A & D & B & A & B & AC & CD & BC & BD\\
\bottomrule
\end{tabular}

\section*{二、必考题}

\noindent
\textbf{11.} (3) 见图线（$l$--$n$ 为一条直线）；\quad(4) 15.35~cm；\quad(5) 127~g

\smallskip

\noindent
\textbf{12.} (1) 大于；\quad(2) $\times 10$；\quad(3) 向上；\quad(4) 400~\si{\ohm}

\smallskip

\noindent
\textbf{13.} (1) $E_{0}=\dfrac{mgd(R_{1}+R_{2})}{qR_{2}}$；\quad(2) $B=\dfrac{mv}{2qd}$；\quad(3) $E'_{\min}=\dfrac{mg}{2q}$

\smallskip

\noindent
\textbf{14.} (1) $k=\sqrt{\dfrac{(1+\lambda)h}{(1-\lambda)H}}$；\quad(2) $F_{0}=\dfrac{2(1-\lambda)(H-h)mg}{h-h_{0}}$

\smallskip

\noindent
(3) $I=m\sqrt{\dfrac{2g(1-\lambda)(H-h)(H^{N+1}-h^{N+1})}{h(H^{N}-h^{N})}}$

\section*{三、选考题}

\noindent
\textbf{15.} (1) ABE；\quad(2)(ⅰ) $p_{1}=1.0\times10^{5}$~Pa；\quad(ⅱ) $F=1$~N

\smallskip

\noindent
\textbf{16.} (1) ABD；\quad(2)(ⅰ) $d=\dfrac{2\sqrt{15}}{5}$~mm；\quad(ⅱ) $x=\dfrac{10\sqrt{3}-12}{15}$~mm
